\begin{center}
\LARGE
Extensible Notations in Mathematica 3.0

\end{center}

\begin{center}
\Large
Jason Harris
\end{center}

\noindent
Date:  July 20th (Saturday)\\
Time: 11:30-12:00\\

\begin{center}
\large
Abstract
\end{center}



Traditionally the use of computer algebra systems has sometimes been
limited in various fields even when applicable because once a problem is
translated from the form in which the user thinks and works on paper into
an ASCII representation suitable for the symbolic computation system most
of the visual intuition is lost. It has long been a goal to allow the use
of notations in which one thinks and works on paper to be used in symbolic
computation.

In this talk I will present a new extension to Mathematica 3.0 that enables
the easy and flexible creation of new notations.  This allows the user to
not only work in such a language but also to write and execute programs in
this language / notation. This is of great benefit to the many users of who
have specialized notations in their field of science.

Several basic examples will be given outlining the capabilities of this
system. I will also discuss the limits of such a system and pose some
questions for future systems... For illustration I will also briefly
present some examples from Logic, Physics and Category Theory.

It is planned that Notation.m will ship as an 'Experimental' package with
Mathematica 3.0.
